% Options for packages loaded elsewhere
\PassOptionsToPackage{unicode}{hyperref}
\PassOptionsToPackage{hyphens}{url}
\PassOptionsToPackage{dvipsnames,svgnames,x11names}{xcolor}
\documentclass[
  11pt,
]{article}
\usepackage{xcolor}
\usepackage[margin=1in]{geometry}
\usepackage{amsmath,amssymb}
\setcounter{secnumdepth}{-\maxdimen} % remove section numbering
\usepackage{iftex}
\ifPDFTeX
  \usepackage[T1]{fontenc}
  \usepackage[utf8]{inputenc}
  \usepackage{textcomp} % provide euro and other symbols
\else % if luatex or xetex
  \usepackage{unicode-math} % this also loads fontspec
  \defaultfontfeatures{Scale=MatchLowercase}
  \defaultfontfeatures[\rmfamily]{Ligatures=TeX,Scale=1}
\fi
\usepackage{lmodern}
\ifPDFTeX\else
  % xetex/luatex font selection
\fi
% Use upquote if available, for straight quotes in verbatim environments
\IfFileExists{upquote.sty}{\usepackage{upquote}}{}
\IfFileExists{microtype.sty}{% use microtype if available
  \usepackage[]{microtype}
  \UseMicrotypeSet[protrusion]{basicmath} % disable protrusion for tt fonts
}{}
\makeatletter
\@ifundefined{KOMAClassName}{% if non-KOMA class
  \IfFileExists{parskip.sty}{%
    \usepackage{parskip}
  }{% else
    \setlength{\parindent}{0pt}
    \setlength{\parskip}{6pt plus 2pt minus 1pt}}
}{% if KOMA class
  \KOMAoptions{parskip=half}}
\makeatother
\usepackage{longtable,booktabs,array}
\newcounter{none} % for unnumbered tables
\usepackage{calc} % for calculating minipage widths
% Correct order of tables after \paragraph or \subparagraph
\usepackage{etoolbox}
\makeatletter
\patchcmd\longtable{\par}{\if@noskipsec\mbox{}\fi\par}{}{}
\makeatother
% Allow footnotes in longtable head/foot
\IfFileExists{footnotehyper.sty}{\usepackage{footnotehyper}}{\usepackage{footnote}}
\makesavenoteenv{longtable}
\setlength{\emergencystretch}{3em} % prevent overfull lines
\providecommand{\tightlist}{%
  \setlength{\itemsep}{0pt}\setlength{\parskip}{0pt}}
\usepackage{bookmark}
\IfFileExists{xurl.sty}{\usepackage{xurl}}{} % add URL line breaks if available
\urlstyle{same}
\hypersetup{
  pdftitle={Normal-ordered neutrino masses from the Koide relation with one negative root{,} and a reason for Q = 2/3},
  pdfauthor={Matthew Gelfman (ORCID 0009-0000-6402-2258); Melissa Behrends-Gelfman (ORCID 0009-0001-6281-1947)},
  colorlinks=true,
  linkcolor={Maroon},
  filecolor={Maroon},
  citecolor={Blue},
  urlcolor={Blue},
  pdfcreator={LaTeX via pandoc}}

\usepackage{amssymb}
\usepackage{newunicodechar}
\newunicodechar{°}{\ensuremath{^\circ}}
\newunicodechar{±}{\ensuremath{\pm}}
\newunicodechar{²}{\ensuremath{^{2}}}
\newunicodechar{³}{\ensuremath{^{3}}}
\newunicodechar{·}{\ensuremath{\cdot}}
\newunicodechar{½}{\ensuremath{\frac{1}{2}}}
\newunicodechar{×}{\ensuremath{\times}}
\newunicodechar{Δ}{\ensuremath{\Delta}}
\newunicodechar{Σ}{\ensuremath{\Sigma}}
\newunicodechar{δ}{\ensuremath{\delta}}
\newunicodechar{θ}{\ensuremath{\theta}}
\newunicodechar{μ}{\ensuremath{\mu}}
\newunicodechar{π}{\ensuremath{\pi}}
\newunicodechar{σ}{\ensuremath{\sigma}}
\newunicodechar{τ}{\ensuremath{\tau}}
\newunicodechar{⁵}{\ensuremath{^{5}}}
\newunicodechar{⁻}{\ensuremath{^{-}}}
\newunicodechar{₁}{\ensuremath{_{1}}}
\newunicodechar{₂}{\ensuremath{_{2}}}
\newunicodechar{₃}{\ensuremath{_{3}}}
\newunicodechar{−}{\ensuremath{-}}
\newunicodechar{√}{\ensuremath{\surd}}
\newunicodechar{≈}{\ensuremath{\approx}}
\title{Normal-ordered neutrino masses from the Koide relation with one
negative root, and a reason for Q = 2/3}
\author{Matthew Gelfman (ORCID 0009-0000-6402-2258) \and Melissa
Behrends-Gelfman (ORCID 0009-0001-6281-1947)}
\date{25 November 2026}

\newunicodechar{⊕}{\ensuremath{\oplus}}
\newunicodechar{¹}{\ensuremath{^{1}}}
\newunicodechar{⁴}{\ensuremath{^{4}}}
\newunicodechar{₀}{\ensuremath{_{0}}}
\newunicodechar{₄}{\ensuremath{_{4}}}
\newunicodechar{₅}{\ensuremath{_{5}}}
\newunicodechar{₆}{\ensuremath{_{6}}}
\newunicodechar{ₖ}{\ensuremath{_{k}}}
\newunicodechar{̂}{}
\newunicodechar{φ}{\ensuremath{\phi}}
\newunicodechar{∮}{\ensuremath{\oint}}
\newunicodechar{‖}{\ensuremath{\Vert}}
\newunicodechar{≡}{\ensuremath{\equiv}}
\newunicodechar{≢}{\ensuremath{\not\equiv}}
\newunicodechar{∎}{\ensuremath{\blacksquare}}
\newunicodechar{𝒯}{\ensuremath{\mathcal{T}}}
\newunicodechar{⁸}{\ensuremath{^{8}}}
\newunicodechar{⁰}{\ensuremath{^{0}}}
\newunicodechar{≤}{\ensuremath{\le}}
\newunicodechar{ε}{\ensuremath{\epsilon}}
\newunicodechar{Γ}{\ensuremath{\Gamma}}
\newunicodechar{∗}{\ensuremath{*}}
\newunicodechar{′}{\ensuremath{^{\prime}}}
\newunicodechar{ᵀ}{\ensuremath{^{T}}}
\newunicodechar{₊}{\ensuremath{_{+}}}
\newunicodechar{₋}{\ensuremath{_{-}}}
\newunicodechar{⁹}{\ensuremath{^{9}}}
\newunicodechar{⁶}{\ensuremath{^{6}}}
\newunicodechar{⁷}{\ensuremath{^{7}}}
\title{Normal-ordered neutrino masses from the Koide relation with one negative
root, and a reason for Q = 2/3}
\author{Matthew Gelfman (ORCID 0009-0000-6402-2258) \\ Melissa Behrends-Gelfman (ORCID 0009-0001-6281-1947)}
\date{Version 1.3, 27 November 2026}
\begin{document}
\maketitle
\hypertarget{abstract}{%
\subsection{Abstract}\label{abstract}}

The charged-lepton masses satisfy Koide's relation Q = (\ensuremath{\Sigma}m)/(\ensuremath{\Sigma}\ensuremath{\surd}m)\ensuremath{^{2}} = 2/3 to about 10\ensuremath{^{-}}\ensuremath{^{5}}. Brannen proposed applying the same relation to neutrinos with one negative square root. Fixed by the two measured mass splittings, it then has a normal-ordered solution:

\begin{itemize}
\tightlist
\item
  with JUNO's 2026 splittings: m\ensuremath{_{1}} = 0.359, m\ensuremath{_{2}} = 8.60, m\ensuremath{_{3}} = 50.09 meV, summing to 59.05 meV;
\item
  with the NuFIT 6.0 global fit (Esteban et al.~2024): m\ensuremath{_{2}} = 8.66 meV and a sum of 59.16--59.36 meV.
\end{itemize}

The inverted-ordering solution sums to about 0.103 eV, which today's cosmological bound disfavours. With TM1 mixing and free Majorana phases, the effective neutrinoless double-beta mass is 1.34--4.02 meV.

We give a reason for Q = 2/3 itself. If the family amplitude matrix is invariant under a 120\ensuremath{^\circ} turn, it is circulant: a\ensuremath{\cdot}1 + b\ensuremath{\cdot}P + b*\ensuremath{\cdot}P\ensuremath{^{2}}. Maximum uncertainty between the self term and the relationship terms then fixes the self share of its weight at \ensuremath{\frac{1}{2}}, which is exactly Q = 2/3. This principle was adopted knowing that Koide's relation holds, so we state it as a choice, not a derivation. Its value is that one statement covers the charged-lepton and neutrino cases together. The predictions can fail through inverted ordering, a sum clearly different from 0.059 eV, or a lightest mass far from 0.36 meV.

\hypertarget{the-relation-and-its-geometry}{%
\subsection{1. The relation and its geometry}\label{the-relation-and-its-geometry}}

Write the square roots of the three masses as the eigenvalues of a Hermitian circulant matrix, \ensuremath{\surd}m\textsubscript{k} = a + 2\textbar b\textbar{} cos(\ensuremath{\delta} + 2\ensuremath{\pi}k/3). The angle between the root vector and the democratic direction (1,1,1) satisfies tan \ensuremath{\theta} = \ensuremath{\sqrt{2}} \textbar b\textbar/a. Koide's Q = 2/3 is \ensuremath{\theta} = 45\ensuremath{^\circ}, that is \textbar b\textbar{} = a/\ensuremath{\sqrt{2}} (Koide 1983; Foot 1994; Brannen 2006).

For the charged leptons, the masses of e and \ensuremath{\mu} predict m\textsubscript{\ensuremath{\tau}} = 1776.969 MeV, against 1776.93 \ensuremath{\pm} 0.09 MeV (+0.4\ensuremath{\sigma}). The twist \ensuremath{\delta} comes out as 0.222222, which matches 2/9 to about one part in a million but is not derived here. Exactly 2/9 together with exactly 45\ensuremath{^\circ} would put m\textsubscript{e} 9.8 ppm low.

\hypertarget{why-the-self-share-is-one-half}{%
\subsection{2. Why the self share is one half}\label{why-the-self-share-is-one-half}}

The self share of the matrix weight is p = 3a\ensuremath{^{2}}/(3a\ensuremath{^{2}} + 6\textbar b\textbar\ensuremath{^{2}}). Every symmetric measure of uncertainty between the two parts, self and relationship, is largest at p = \ensuremath{\frac{1}{2}} (Shannon 1948). And p = \ensuremath{\frac{1}{2}} is exactly \textbar b\textbar\ensuremath{^{2}} = a\ensuremath{^{2}}/2, the 45\ensuremath{^\circ} condition. The measured charged leptons give p = 0.50000.

For quarks, p = 0.456 (down-type) and 0.393 (up-type), and no choice of signs reaches \ensuremath{\frac{1}{2}}. The principle is therefore stated for colour-neutral families only. Running from a high scale should spoil the lepton relation by about 0.35\% (Sumino 2009). It holds for pole masses, and we do not explain why.

\hypertarget{neutrinos}{%
\subsection{3. Neutrinos}\label{neutrinos}}

With one negative root, the conditions Q = 2/3, m\ensuremath{_{2}}\ensuremath{^{2}} \ensuremath{-} m\ensuremath{_{1}}\ensuremath{^{2}} = \ensuremath{\Delta}m\ensuremath{^{2}}\ensuremath{_{2}}\ensuremath{_{1}} and m\ensuremath{_{3}}\ensuremath{^{2}} \ensuremath{-} m\ensuremath{_{1}}\ensuremath{^{2}} = \ensuremath{\Delta}m\ensuremath{^{2}}\ensuremath{_{3}}\ensuremath{_{1}} have a unique normal-ordered solution:

\begin{longtable}[]{@{}
  >{\raggedright\arraybackslash}p{(\columnwidth - 8\tabcolsep) * \real{0.2000}}
  >{\raggedright\arraybackslash}p{(\columnwidth - 8\tabcolsep) * \real{0.2000}}
  >{\raggedright\arraybackslash}p{(\columnwidth - 8\tabcolsep) * \real{0.2000}}
  >{\raggedright\arraybackslash}p{(\columnwidth - 8\tabcolsep) * \real{0.2000}}
  >{\raggedright\arraybackslash}p{(\columnwidth - 8\tabcolsep) * \real{0.2000}}@{}}
\toprule\noalign{}
\begin{minipage}[b]{\linewidth}\raggedright
Splittings used
\end{minipage} & \begin{minipage}[b]{\linewidth}\raggedright
m\ensuremath{_{1}} (meV)
\end{minipage} & \begin{minipage}[b]{\linewidth}\raggedright
m\ensuremath{_{2}} (meV)
\end{minipage} & \begin{minipage}[b]{\linewidth}\raggedright
m\ensuremath{_{3}} (meV)
\end{minipage} & \begin{minipage}[b]{\linewidth}\raggedright
\ensuremath{\Sigma} (meV)
\end{minipage} \\
\midrule\noalign{}
\endhead
\bottomrule\noalign{}
\endlastfoot
JUNO 2026 (\ensuremath{\Delta}m\ensuremath{^{2}}\ensuremath{_{2}}\ensuremath{_{1}} = 7.388 \ensuremath{\times} 10\ensuremath{^{-}}\ensuremath{^{5}}, \ensuremath{\Delta}m\ensuremath{^{2}}\ensuremath{_{3}}\ensuremath{_{1}} = 2.509 \ensuremath{\times} 10\ensuremath{^{-}}\ensuremath{^{3}} eV\ensuremath{^{2}}) & 0.359 & 8.60 & 50.09 & 59.05 \\
NuFIT 6.0, without SK atmospheric data & 0.364 & 8.66 & 50.13 & 59.16 \\
NuFIT 6.0, with SK atmospheric data & 0.362 & 8.66 & 50.34 & 59.36 \\
\end{longtable}

The sum sits just above the normal-ordering minimum (58.8 meV), so the most distinctive number is the lightest mass. The inverted solution (\ensuremath{\Sigma} \ensuremath{\approx} 0.103 eV) is disfavoured by the cosmological bound, not by the relation.

With TM1 mixing (Xing and Zhou 2014; Albright and Rodejohann 2009) and free Majorana phases, the effective double-beta mass is 1.34--4.02 meV, or 1.82 meV with zero phases and \ensuremath{\delta} = 270\ensuremath{^\circ}. This is below the KamLAND-Zen limit of 28--122 meV (KamLAND-Zen Collaboration 2024) and below planned reach. The masses themselves do not depend on TM1.

\hypertarget{tests}{%
\subsection{4. Tests}\label{tests}}

\begin{itemize}
\tightlist
\item
  \textbf{Ordering:} after its first measurement of the solar parameters (JUNO Collaboration 2025), JUNO favours normal ordering at 2.1\ensuremath{\sigma} (JUNO Collaboration 2026) and expects 3\ensuremath{\sigma} after about six and a half years.
\item
  \textbf{Sum:} DESI, Euclid and next-generation CMB experiments approach the needed sensitivity.
\item
  \textbf{Majorana nature:} proof that neutrinos are Dirac particles would remove the double-beta prediction.
\end{itemize}

Brannen's proposal and the present statement give the same numbers. This test can falsify both but cannot separate them. What this note adds is one principle behind both the charged-lepton and neutrino cases, and an up-to-date evaluation.

\hypertarget{data-and-code}{%
\subsection{Data and code}\label{data-and-code}}

\texttt{analysis/recalibrate\_to\_2026\_values\_2026-11-18.py} and \texttt{analysis/derive\_45\_and\_sqrt2\_2026-11-18.py} in the TRPG code archive, doi:10.5281/zenodo.24329239. This result is one component of the framework in doi:10.5281/zenodo.24330763.

\hypertarget{acknowledgements}{%
\subsection{Acknowledgements}\label{acknowledgements}}

Due to medical necessity, AI language-model assistants were used, under the authors' direction, to draft and edit text from the authors' ideas, to carry out and check calculations, and to search the literature; every result was checked by rerunning the published code. The authors reviewed all content and take full responsibility for it.

This paper is licensed under CC BY 4.0; its scripts and data are under the MIT License.

\hypertarget{references}{%
\subsection{References}\label{references}}

\begin{itemize}
\tightlist
\item
  Albright, C. H. and Rodejohann, W. (2009). Comparing trimaximal mixing and its variants with deviations from tri-bimaximal mixing. \emph{Eur. Phys. J. C} 62, 599.
\item
  Brannen, C. A. (2006). The lepton masses. Online note, brannenworks.com.
\item
  Esteban, I., Gonzalez-Garcia, M. C., Maltoni, M., Martinez-Soler, I., Pinheiro, J. P. and Schwetz, T. (2024). NuFit-6.0: updated global analysis of three-flavor neutrino oscillations. \emph{JHEP} 12, 216.
\item
  Foot, R. (1994). A note on Koide's lepton mass relation. arXiv:hep-ph/9402242.
\item
  JUNO Collaboration (2025). Measurement of reactor neutrino oscillation with the first JUNO data. \emph{Nature} 654, 343 (2026); arXiv:2511.14593.
\item
  JUNO Collaboration (2026). Measurement of neutrino oscillation parameters at JUNO and indication of the neutrino mass ordering. arXiv:2609.37895.
\item
  KamLAND-Zen Collaboration (2024). Search for Majorana neutrinos with the complete KamLAND-Zen dataset. \emph{Phys. Rev.~Lett.} 135, 262501 (2025); arXiv:2406.11438.
\item
  Koide, Y. (1983). A fermion--boson composite model of quarks and leptons. \emph{Phys. Lett. B} 120, 161.
\item
  Shannon, C. E. (1948). A mathematical theory of communication. \emph{Bell Syst. Tech. J.} 27, 379.
\item
  Sumino, Y. (2009). Family gauge symmetry and Koide's mass formula. \emph{Phys. Lett. B} 671, 477.
\item
  Xing, Z.-z. and Zhou, S. (2014). A partial \ensuremath{\mu}--\ensuremath{\tau} symmetry and its prediction for leptonic CP violation. \emph{Phys. Lett. B} 737, 196.
\end{itemize}

\end{document}
